% Generated draft: requires mathematical and source review.
\documentclass[11pt]{article}
\usepackage[margin=25mm]{geometry}
\usepackage{fontspec}
\setmainfont{Latin Modern Roman}
\usepackage{amsmath,amssymb,mathtools,unicode-math}
\setmathfont{Latin Modern Math}
\usepackage{enumitem,needspace,xurl,hyperref}
\setlength{\parindent}{0pt}
\setlength{\parskip}{5pt}
\setlength{\emergencystretch}{3em}
\newcommand{\N}{\mathbb{N}}
\newcommand{\Z}{\mathbb{Z}}
\newcommand{\Q}{\mathbb{Q}}
\newcommand{\R}{\mathbb{R}}
\newcommand{\C}{\mathbb{C}}
\newcommand{\F}{\mathbb{F}}
\newcommand{\A}{\mathbb{A}}
\newcommand{\PP}{\mathbb{P}}
\newcommand{\E}{\mathbb{E}}
\newcommand{\eps}{\varepsilon}
\newcommand{\dd}{\,\mathrm{d}}
\begin{document}
\section*{UnsolvedMath ID: 700715}
\textbf{Source problem code:} catalogue entry
\subsection{Definitions and mathematical statement}
Let \(p\) be a prime and let \(n\geq 1\) be an integer. A division ring \(D\) is a ring with identity in which every nonzero element is invertible; commutativity is not assumed. Write
\[
T_n(D)=\{(a_{ij})\in\operatorname{GL}_n(D):a_{ij}=0\text{ whenever }i>j\}
\]
for its group of invertible upper triangular matrices. A faithful representation of a group \(G\) in \(T_n(D)\) is an injective group homomorphism \(\rho:G\to T_n(D)\).

Under the torsion convention, a \(p\)-group is a group in which every element has finite order equal to a power of \(p\). The source does not explicitly specify whether only finite \(p\)-groups are intended; this scope qualification requires review.

For each prime \(p\) and positive integer \(n\), characterize the \(p\)-groups \(G\) for which there exist a division ring \(D\) of characteristic \(p\) and an injective homomorphism
\[
\rho:G\hookrightarrow T_n(D).
\]
Here the division ring may depend on \(G\); the formulation treats \(n\) as a prescribed dimension parameter.

The supplied source states that the case \(p\geq n\) was characterized by B. A. F. Wehrfritz in the 1987 reference below. It also reports that the corresponding characteristic-zero question, written there as \(p=0\), was solved by B. Hartley and P. Menal, with a second proof in Wehrfritz's reference. Since zero is not a prime and the source does not specify the group class for that corresponding question, this remark is retained as historical context rather than a characteristic-zero extension of the displayed problem.
\subsection{Short English statement}
For a prime \(p\) and a prescribed dimension \(n\), characterize the \(p\)-groups that embed in the invertible triangular matrices over some division ring of characteristic \(p\). The source reports a characterization when \(p\geq n\). Whether infinite \(p\)-groups are included requires clarification.
\subsection{Sources}
\begin{itemize}
\item Supplied local UnsolvedMath registry record, permanent ID \(700715\), exact source code \texttt{catalogue entry}, title \emph{Kourovka Notebook Problem 11.22}. The registry labels the problem \texttt{open}; this is supplied metadata, not an independently verified assessment.
\item \emph{The Kourovka Notebook (Unsolved Problems in Group Theory)}, original Problem 11.22, 11th issue (1990); supplied edition reference: 21st issue edition (2026), \texttt{arXiv:1401.0300v48}. Proposed by B. A. F. Wehrfritz. The supplied collection description names E. I. Khukhro and V. D. Mazurov as editors.
\item B. A. F. Wehrfritz, \emph{Bull. London Math. Soc.} \textbf{19} (1987), 320--324. The supplied statement attributes the solution for \(p\geq n\), and a second proof for the corresponding characteristic-zero question, to this reference. No article title is supplied.
\item B. Hartley and P. Menal, \emph{Bull. London Math. Soc.} \textbf{15} (1983), 378--383. The supplied statement attributes a solution of the corresponding characteristic-zero question to this reference. No article title is supplied.
\end{itemize}
\end{document}
